\chapter{Introduction}
\label{ch:intro}

\section{The question}

Part~I \cite{PartI}, cited here by its printed numbers, gives two readings of a vector $X\in\R^{3,3}$. They assign different physical
quantities and different Lorentz actions to the same six real components. In the displacement reading, $X$ describes an interval with three
temporal and three spatial components, and a chosen four-plane carries the interior Lorentz algebra $\lint$. In the field reading,
$X=(\bB;\bE)$ represents an electromagnetic field strength, and the Lorentz algebra is $\lG$ (\PI{Thm~3.3}). The two Lorentz actions are not
similar (\PI{Prop.~3.6}), and the readings differ in units, in the quadratic form that is the energy, and in what is configuration and what is
momentum (\PI{Thm~3.15}).

A common operator motivates the comparison. The interior dynamics of \PI{Chapter~4} uses the shift $\partial_\tau+\kappa G$, with $G$ the
binding root (\PI{Thm~4.8}) and $\kappa=c/\ell$ a rate that nothing in the algebra fixes (\PI{Prop.~4.49}). In the field reading the
source-free Maxwell evolution takes the form $\partial_tX=-G\,\curl X$, together with the two Gauss constraints (\PI{Prop.~3.10}). On a
transverse curl eigenfield with eigenvalue $\kappa$ the two evolution laws have the same form once their evolution parameters are identified.
Part~I does not connect the two: \PI{Section~8.6} records that the link between the field-strength reading and the internal connection is not
made, and \PI{Section~13.4} that the two readings meet as base and fibre.

We ask how far this correspondence extends to dynamics and to the internal action on spinors. We first study two field models: a
differential shift involving the curl, and a constant-rate model with a massive two-form sector. The growth of particular tree amplitudes in
the minimally coupled massive model motivates an $\so3$ Yang--Mills--Higgs completion. We then give a transport prescription for the temporal
Clifford generators, and we specify which rotations are gauged and which mass direction is used. Its consequences can be tested separately
on each spinor module of the tower $\Cl(3n,3n)$ of Part~I. The central example is the second rung, where the full module carries a chiral
Pati--Salam family.

\section{Summary of results}

\paragraph{The field reading of the interior dynamics.} We replace $\kappa G$ by $G\circ\curl$, with which it coincides on the corresponding
curl eigenspace. With the signs of \cref{eq:sources}, $A_+X=-(\bk;\bj)$: the momentum $\pi=-m(\bk;\bj)$ of \PI{Chapter~4} contains the
spatial magnetic and electric currents, and the force law is a mirror of Maxwell's equations. The charge densities complete $\pi$ into an
eight-component object $\Pi$ that carries both four-currents, and the split-form Lagrangian completed by them agrees with the free kinetic
term of an $\R^{3,3}$-valued field up to a divergence (\cref{ch:curl}). Read as a mass instead, the shift gives a theory that is the sum of
two halves; the half on the temporal triple is a massive two-form with positive canonical energy, and the binding root maps it to a mirror of
negative canonical energy (\cref{ch:mass}). Coupled to a connection at a constant rate the temporal half keeps its constraint structure and a
positive canonical matter energy (\cref{ch:bcoupling}), but its tree amplitudes grow with energy, in transverse and in longitudinal channels
(\cref{ch:quantum}). An $\so3$ Yang--Mills--Higgs completion, in which the connection and the massive pair are components of one field with a
Higgs vector of variable length, cancels the leading growth in the channels and fixed nonexceptional kinematics studied. Its field equations are the Yang--Mills equations and Bianchi identities in electric and magnetic components: the Maxwell system
of Part~I with covariant derivatives and a current made from the Higgs vector (\cref{ch:connection}).

\paragraph{Clifford transport.} The hypothesis of this Part has four parts, stated precisely in \cref{def:hyp} and referred to
throughout as (A1) to (A4). (A1) \emph{Transport}: the Clifford algebra of the temporal directions acts on the spinors through a map
$\sigma_n$ that replaces the clock by the volume element of the full algebra. (A2) \emph{Yang--Mills}: the gauge field of the temporal
directions takes values in the image $\gn$ of their rotations. (A3) \emph{Mass}: the generator of mass is the image $G_n=\sigma_n(T_n)$ of
the temporal volume; where this applies, it prescribes a vacuum direction of the mass field, without determining the full vacuum. (A4) \emph{The spatial directions}: the spatial triple of the four-plane
carries no gauge field and the mass is a Lorentz-scalar field, while the spatial directions of the other constituents carry the gauge
algebra $\mathfrak s_n$. The transport map $\sigma_n$ replaces the clock $E_1$ of the four-plane by the volume element $\Gamma_n$ of the
full algebra $\Cl(3n,3n)$, with $\Gamma_n^2=\id$, in the whole Clifford algebra of the $3n$ temporal directions (\cref{def:hyp}); it is not the
volume $\Omega$ of the four-plane, for which $\Omega^2=-\id$. At one constituent the transported temporal algebra is the whole commutant of the
interior Lorentz algebra; inside that commutant the admissible symmetric involution is fixed up to sign, and the chosen orientation
selects $\Gamma$. The transported rotations are the algebra of
Pauli and G\"ursey, and the image of the temporal volume $T$, the generator of the Dirac rest rotation, is $-\Omega$, which acts on the
embedded field strength by left multiplication as the binding root (\cref{thm:clifford1,prop:unique,thm:massgen}). The Lorentz-invariant mass
kernels are then exactly the transported Higgs vectors and Yang--Mills generators, each multiplied by the volume of the spatial triple
(\cref{thm:massfield}). At rung $n$ the commutant of the Lorentz algebra is $\Cl(3n,3n-3)$, of dimension $2^{6n-3}$, which contains the
transported temporal algebra, of dimension $2^{3n}$. The transported rotations form $\so{3n}$, and the transported temporal volume is a
rotation for $n\equiv1,2$ and an involution for $n\equiv3,0$ modulo four (\cref{thm:commutant,thm:everyrung,thm:massrung}).

\paragraph{One Pati--Salam family.} At the second rung the transported rotations form $\so6\cong\su4$, and the spatial directions of the
second constituent carry $\su2_L\oplus\su2_R$. With the spectator Lorentz action and (A2), (A4), the left-handed Weyl multiplicity space of
$\R^{64}$ carries exactly $(\mathbf4,\mathbf2,\mathbf1)\oplus(\bar{\mathbf4},\mathbf1,\mathbf2)$, one family of the Standard Model with a
right-handed neutrino, and no projector is chosen (\cref{thm:rung2,thm:family}). Its perturbative gauge anomalies cancel, and each weak factor
has four Weyl doublets, which passes the usual test of the $SU(2)$ global anomaly on spin spacetimes (\cref{thm:ps}). The Lorentz-invariant
mass field decomposes into the scalar representations of left--right symmetric Pati--Salam models. For the chosen Standard-Model embedding
the invariant kernels form one complex line, a right-handed Majorana direction $\Phi_0$. The bidoublet direction $G_2\,\omega_S$ prescribed by
the generator of mass has the largest gauge stabiliser among the nonzero kernels. Both values are needed: together they leave
$\su3\oplus\uu1_Q$, which fixes the generator $Q$ in this specified broken phase, and with equal couplings $\sin^2\theta_W=\tfrac38$ at their
matching scale (\cref{thm:psstage,prop:ewdir,thm:vectors}). The electroweak numbers that Part~I found at the third rung, on a block chosen by a
projector, recur at the second rung on the whole module, with the same trace conventions and vacuum values.

\paragraph{Masses and the obstruction to the sum rule.} The generator of mass alone gives the family a single Dirac entry. In the restricted
real ansatz of the neutral Yang--Mills values the signed Dirac entries are $z_f=\kappa-X(f_L)$, the physical Dirac masses being $|y\,z_f|$,
and they obey one linear sum rule (\cref{thm:split}). No symmetry protects it: the excluded direction $D$ is fixed by the stabiliser of the
vacuum, and the gradient of the invariant trace quartic $\Tr((mm^\top)^2)$, computed in closed form, has a component along it,
$-\half v_R^2\kappa$ at $v_R\Phi_0+\kappa\,G_2\,\omega_S$ (\cref{thm:potential}); the value it induces depends on the couplings and the
restoring curvatures. With an independent coefficient of $D$, all four signed entries of a family can be fitted (\cref{thm:splitgen}), as an
illustration does by fitting the nine charged masses and chosen neutrino inputs (\cref{ex:observed}). Holding the interaction of the second rung fixed, the third rung gives eight
chiral copies of the family; on restriction to the embedded Pati--Salam algebra built with $\Gamma_3$ it gives four families and four
mirrors, which the full third-rung gauge action mixes (\cref{thm:families}).

\section{Rungs as interactions}
\label{sec:interactions}

Part~I treats a rung of the tower as an arena: it asks what the spinor module of one rung can hold, and at the third rung it finds a
Pati--Salam family on a block chosen by a projector, with the weak directions among the temporal generators and colour among the spatial
ones (\PI{Chapter~7}). Here a rung specifies a representation space and a candidate gauge action on it: the algebras that the hypothesis
gauges there and the mass that its generator of mass carries (\cref{fig:rungs}). A requirement that belongs to one interaction is not imposed
on another. Chirality is required of the weak interaction, which the second rung carries, and the Higgs vector at the clock belongs to the
first rung. These considerations motivate the rung-dependent prescription in (A3): at the second rung it supplies the bidoublet
direction, and the additional Majorana value is specified independently; the rule (A4) is kept at the third rung (\cref{def:hyp}). We compare the
consequences of applying the same prescription at successive rungs. Whether the corresponding gauge fields are related by dynamics is a
separate question; in particular, \cref{prop:rung1charge} distinguishes the first-rung generator $\half J$ from the electromagnetic generator
in the broken phase of the second rung.

\begin{figure}[t]
\centering
\begin{tikzpicture}[font=\footnotesize,>=Stealth,
  rung/.style={draw,rounded corners=3pt,align=left,text width=4.5cm,minimum height=3.8cm,inner sep=5pt,fill=blue!3,anchor=north west}]
\node[rung] (r1) at (0,0) {\textbf{Rung one}: $\Cl(3,3)$ on $\R^8$\\[3pt]
  gauged: $\su2$ (Pauli--G\"ursey)\\
  generator: $-\Omega$, a rotation\\
  vacuum: the clock, kernel $\kappa\,T$\\
  unbroken: $\uu1_J$\\
  matter: one Weyl doublet\\
  \quad(global $SU(2)$ anomaly)\\
  \phantom{x}};
\node[rung] (r2) at (5.45,0) {\textbf{Rung two}: $\Cl(6,6)$ on $\R^{64}$\\[3pt]
  gauged: $\su4\oplus\su2_L\oplus\su2_R$\\
  generator: $-\Omega\,\omega_\cN$, a rotation\\
  vacuum: $v_R\Phi_0+\kappa\,G_2\,\omega_S$\\
  \quad(chosen, $v_R\kappa\ne0$)\\
  unbroken: $\su3\oplus\uu1_Q$\\
  matter: one chiral family\\
  \quad(perturbative anomalies cancel;\\
  \quad even number of doublets\\
  \quad per weak factor)};
\node[rung] (r3) at (10.9,0) {\textbf{Rung three}: $\Cl(9,9)$ on $\R^{512}$\\[3pt]
  gauged: $\so9\oplus\so8$\\
  generator: $G_3$, an involution\\
  vacuum: (A3) selects none\\
  an invariant mass kernel\\
  matter: $(\mathbf{16},\mathbf8)$\\
  \quad(real; admits an invariant\\
  \quad mass)\\
  \phantom{x}};
\draw[->,thick] (r1.east) -- node[above,font=\scriptsize]{$+1$} (r2.west);
\draw[->,thick] (r2.east) -- node[above,font=\scriptsize]{$+1$} (r3.west);
\end{tikzpicture}
\caption{The first three rungs of the tower read as candidate gauge actions on successive modules, under the hypothesis of \cref{def:hyp};
``generator'' is the generator of mass $G_n$ of (A3), and each arrow adds a constituent. The second-rung vacuum is a chosen two-stage
background: (A3) supplies the bidoublet direction, and the right Majorana value is an additional choice. The entries are established in
\cref{ch:hyp} (rung one), \cref{ch:rungs,ch:ew} (rung two) and \cref{ch:rungs} (rung three).}
\label{fig:rungs}
\end{figure}

\section{Conventions}
\label{sec:conv}

The conventions are those of Part~I. $\Cl(3,3)\cong M_8(\R)$ acts on $S=\R^8$ through the generators $e_1,\dots,e_6$ of \PI{eq.~(1.1)},
with $\eta=\mathrm{diag}(1,1,1,-1,-1,-1)$; $e_1,e_2,e_3$ are symmetric matrices and span the temporal triple $V^+$, and $e_4,e_5,e_6$ are
antisymmetric and span the spatial triple $V^-$. In $\Cl(p,q)$ the first index counts generators of square $+\id$, so that
$\Cl(3,0)\cong M_2(\C)$, $\Cl(6,0)\cong M_4(\mathbb H)$ and $\Cl(0,6)\cong M_8(\R)$. The clock is a unit vector $\clock$ of $V^+$
(\PI{Assumption~2.16}), taken as $e_1$; the four-plane is $E=\langle e_1;e_4,e_5,e_6\rangle$ with $\gamma^0=e_1$ and $\gamma^a=e_{a+3}$, and
the surplus pair is $(e_2,e_3)$. We write
\[
J=e_2e_3,\qquad T=e_1e_2e_3,\qquad \omega_S=e_4e_5e_6,\qquad \Omega=e_1e_4e_5e_6=\gamma^0\gamma^1\gamma^2\gamma^3,\qquad \Gamma=e_1e_2\cdots e_6 .
\]
The orientation of the four-plane is the one fixed by $\Omega$; with it $\Omega$ acts on the two-forms as minus the Hodge star
(\PI{Prop.~3.7}), and this orientation is used wherever the sign of the duality matters. The interior Lorentz algebra $\lint$ is spanned by
the rotations $\half\gamma^a\gamma^b$ and the boosts $\half\gamma^0\gamma^a$; the kinetic forms are $K^0=-\id$ and $K^a=-\gamma^0\gamma^a$, and
$T=\gamma^0J$ is the Dirac mass form.

In the field reading $X=\sum_iB_ie_i+E_ie_{i+3}$, written $X=(\bB;\bE)$, with $\eta(X,X)=|\bB|^2-|\bE|^2$ and $\bB=-\Btext$, where $\Btext$
is the textbook magnetic field. The binding root is $G(\bB;\bE)=(-\bE;\bB)$; the complex coordinate is $z=\bB+i\bE$, on which $G$ is
multiplication by $i$; and $\lG=\mathrm{span}\{L_i,K_i=GL_i\}$. Sources are written with both kinds of current,
\begin{equation}
\label{eq:sources}
\partial_t\bB=\curl\bE-\bk,\qquad \partial_t\bE=-\curl\bB-\bj,\qquad \rho_e=\dive\bE,\qquad \rho_m=\dive\bB ,
\end{equation}
with the four-currents $J_e=(\rho_e,\bj)$ and $J_m=(\rho_m,\bk)$ and the Minkowski squares $J\cdot J=\rho^2-|\text{current}|^2$. The
undivided energy is $E_6=\half(|\bB|^2+|\bE|^2)$, and Maxwell's Lagrangian at the field point is $L_M=-\half\eta(X,X)$ (\PI{Prop.~3.20}).
Throughout $c=1$; in \cref{ch:quantum,ch:connection} also $\hbar=1$.

Three placements of the pair $(\bB;\bE)$ occur in Part~I and are used here. In the frame vector $X\in V$, $\bB$ sits on the temporal triple
and $\bE$ on the spatial one (\PI{Thm~3.3}); this is the placement of \cref{ch:curl,ch:mass,ch:bcoupling,ch:quantum,ch:connection}. In the
two-form $F_X$ on the four-plane, $\bE$ sits on the planes that contain the clock, the boosts $\gamma^0\gamma^i$, and $\bB$ on the purely
spatial planes, as $F_{0i}=E_i$ and $F_{jk}=\epsilon_{ijk}B_i$ require; this is the placement of \cref{thm:massgen}. In the spinor module,
$\bE$ sits on the three octants of Fock weight one and $\bB$ on the three of weight two (\PI{Introduction and Section~3.5}); the sign triples
that label these octants are eigenvalues of the pole operators $e_ie_{i+3}$, not metric signatures. The frame placement carries the signs of
the tensor convention: $\eta(X,X)=\half F_{\mu\nu}F^{\mu\nu}$ in the signature $(+,-,-,-)$, and in rationalised units with $c=1$ the free
Maxwell density is $-\tfrac14F_{\mu\nu}F^{\mu\nu}=-\half\eta(X,X)$; Jackson, eq.~(12.85), and Landau and Lifshitz, eq.~(27.4), write the
corresponding expressions in Gaussian units \cite{Jackson,LandauLifshitz2}. The opposite sign, $|\bE|^2-|\bB|^2$, is the real part of the square of the
Riemann--Silberstein vector $\bE+i\bB$ \cite{Silberstein1907,BialynickiBirula13} and of the bivector $\bE+I\bB$ of the spacetime algebra
\cite{Hestenes66}; it is the sign of the second placement.

At rung $n$ the algebra is $\Cl(3n,3n)$ on $\R^{8^n}$, with the generators
$E_{6(k-1)+a}=\omega^{\otimes(k-1)}\otimes e_a\otimes\id^{\otimes(n-k)}$, $\omega=e_1\cdots e_6$, for the constituent $k$ and $a=1,\dots,6$.
The four-plane, $\Omega$ and $\lint$ are those of the first constituent: this is the spectator action of \PI{Section~5.2}, as used at the
second rung in \PI{Section~8.5.3}. The $3n-1$ positive generators outside the four-plane are called internal, and their set is $P$: the
surplus pair $E_2,E_3$ and the temporal triples of the other constituents. The set of the $3n-3$ spatial generators of the other constituents
is $\cN$, and $\omega_\cN$ is their product. As in \PI{Thm~9.62}, $T_n$ is the product of the $3n$ temporal generators in increasing order,
$S_n$ that of the $3n$ spatial ones, and $\Gamma_n$ is the volume element of $\Cl(3n,3n)$; $T_1=T$, $S_1=\omega_S$ and $\Gamma_1=\Gamma$. The
index $n$ is dropped where the rung is clear. Two charges must be kept apart: the charge $\half J$ of the first rung (\cref{sec:rung1int})
and the electric charge $Q=T_{3L}+T_{3R}+\half(B\!-\!L)$ of the second (\cref{prop:rung1charge}); both appear in \cref{fig:rungs}.

\section{Scope and evidence}
\label{sec:scope}

This Part studies the consequences of the transport and model assumptions (A1) to (A4), with the spectator Lorentz action, together with
the field reading of Part~I. It makes no empirical prediction. The value $\sin^2\theta_W=\tfrac38$ is a matching value under the assumption
of a single coupling at an unspecified scale; the masses of \cref{ex:observed} are fitted, not derived. The scattering and mass-spectrum
calculations are at tree level; \cref{ch:rungs,ch:ew} also evaluate anomaly coefficients. No statement of Part~I is changed: where a result
here bears on one, the bearing is stated, and the statement of Part~I stands as printed.

Each statement is supported in one of three ways: by a proof in the text; by a scripted computation, named in a line \emph{Checks} that
closes the statement or its proof; or by a citation, when the result is standard. Interpretations are stated as such in the prose, and the
checks do not test them. A check establishes the identity or the number that it states and nothing more. Symbolic identities and matrix
identities with integer or dyadic entries are exact; rank decisions from floating-point singular values, sampled amplitudes and numerical
searches hold with the stated tolerances and on the stated domains; statements for general $n$ are proved in the text and checked at
$n=1,2,3$ (\cref{app:verification}).

\section{Relation to prior work}
\label{sec:prior}

Several programmes read Standard-Model structure off a Clifford or division algebra: the gauge symmetries of the Standard Model from the
rotations of the spinors of the Clifford algebra $C\ell_7$ \cite{TraylingBaylis01}, leptons and quarks from the ideals of the complex Clifford
algebra $\C\ell_6$ \cite{Stoica18}, and constructions from the normed division algebras \cite{Furey,Gresnigt}. The representation theory of
one family, with the branching of the spinor of $\so{10}$ through $\so6\oplus\so4$ to the Pati--Salam algebra, is reviewed in
\cite{BaezHuerta10}. The spectral programme obtains a Pati--Salam model with sixteen fermions per family when the order-one condition is
dropped \cite{ChamseddineConnesVanSuijlekom13}; \PI{Section~14.4} compares it with the constructions of Part~I. The construction closest to
\cref{ch:rungs} is Krasnov's \cite{Krasnov18}. He realises the spinors of $SO(7,7)$ on the forms of a null seven-plane, finds the Pati--Salam
algebra as the compact part $\so4\oplus\so6$ of the commutant $\so{4,6}$ of a Lorentz subalgebra, and places one family with a right-handed
neutrino on a module of real dimension $64$; the Lorentz subalgebra is selected by a choice that is not made dynamical. \PI{Section~14.3}
shows that his module and the chiral block of Part~I are equivalent as representations of $\so{1,3}\oplus\so4\oplus\so6$, once the outer
automorphisms are fixed compatibly on both sides. The family of \cref{thm:rung2} has the same content, on the whole spinor module of the
second rung, and the same argument applies to it.

The matter content is therefore shared prior art. What this Part adds is the route to it and what follows from it: the transport map and its
commutant theorems at one constituent and at every rung (\cref{thm:clifford1,thm:massgen,thm:commutant,thm:everyrung}), which relate the
temporal rotations of the displacement reading to an internal gauge algebra; the family on the whole module of the second rung without a
projector, under (A1) to (A4) (\cref{thm:rung2,thm:ps}); the analysis of the mass field, with a Majorana line that the transported algebra
does not reach (\cref{thm:psstage,prop:notransport}); and the explicit obstruction to the mass relation of the restricted ansatz
(\cref{thm:potential}). The field dynamics of \cref{ch:curl,ch:mass,ch:bcoupling,ch:quantum,ch:connection} continues \PI{Chapter~4}; the
results it inherits, from Maxwell's and Proca's equations to tree unitarity and the model of Georgi and Glashow, are cited where they are
used.

\section{Organisation}

\Cref{ch:curl,ch:mass,ch:bcoupling,ch:quantum,ch:connection} read the interior dynamics of Part~I on fields: the shift as the curl, the shift
as a mass, the coupling to a connection, the quantum status of that coupling, and its completion. \Cref{ch:hyp} states the hypothesis and
works it out at one constituent. \Cref{ch:rungs} climbs the tower, with the second rung and the third. \Cref{ch:ew} realises the electroweak
stage of the second rung and the masses of the family, with the obstruction to the sum rule. \Cref{ch:conclusions} collects the assumptions,
the obstructions and the open problems, and \cref{app:verification} describes the verification and the code. \Cref{sec:prior} places the
construction among earlier ones.

\section{Use of AI assistance}
\label{sec:ai}

This Part continues Part~I and was produced by the same method (\PI{Section~I.3 and Chapter~15}). It was developed in an iterative process
with AI systems: Claude models assisted with candidate statements, exploratory proofs, alternative constructions, the verification scripts,
calculations and the drafting of the text, and ChatGPT and Perplexity with discussion and with referee-style criticism of successive drafts. The questions,
the hypothesis and the decisions on what is stated were not delegated. Every candidate result was examined, and then proved in the text,
reduced to a scripted check, referred to the literature, or discarded; objections raised in the reviews were checked by independent
computation before the text was corrected. No particular model version was in use throughout. AI-assisted review is not peer review, and
agreement between models is not proof: each statement of this Part carries its own evidence, as described in \cref{sec:scope}. I take full
responsibility for the content of this Part and for its conclusions.
